\part{Going further}

% =====================================================================
\chapter{Writing your own kinetics}
\label{ch:ownkinetics}

\section{The thing to understand first}

\begin{cbwarn}
\textbf{\color{cbRed}Rate laws are compiled in, not read from the input
file.}\quad
\file{defineKinetics.hh} and \file{defineAbioticKinetics.hh} are
\code{\#include}d into the executable. Changing your chemistry means editing a
C\texttt{++} header and rebuilding. Two problems with different chemistry need
two binaries.

This is the single most surprising thing about the code, and it is why every
example folder carries its own pair of headers and why
\file{runAllExamples.sh} rebuilds between cases.
\end{cbwarn}

\section{The two hooks}

\begin{center}
\begin{tabular}{@{}p{4.6cm}p{10.2cm}@{}}
\toprule
\code{defineRxnKinetics} & biotic. Gets the biomass vector as well as the
concentrations, and writes growth rates. Runs when \tg{enable\_kinetics} is
true, for organisms whose \tg{reaction\_type} includes kinetics.\\
\code{defineAbioticRxnKinetics} & abiotic. Concentrations only. Runs when
\tg{enable\_abiotic\_kinetics} is true.\\
\code{defineDissolutionRate} & the dissolution hook, in the abiotic file. Runs
only when \tg{dissolution} is enabled.\\
\bottomrule
\end{tabular}
\end{center}

\section{Anatomy of a biotic rate law}

\begin{lstlisting}[style=cpp]
void defineRxnKinetics(
    std::vector<double>  B,      // IN : biomass here, gDW/L, one per microbe
    std::vector<double>  C,      // IN : concentrations here, XML order
    std::vector<double>& subsR,  // OUT: dC/dt per substrate, per SECOND
    std::vector<double>& bioR,   // OUT: dB/dt per microbe,   per SECOND
    plb::plint           mask)   // IN : material number of this voxel
{
    for (std::size_t i = 0; i < subsR.size(); ++i) subsR[i] = 0.0;
    for (std::size_t i = 0; i < bioR.size();  ++i) bioR[i]  = 0.0;
    if (C.size() < 2 || subsR.size() < 2) return;   // guard: your chemistry size
    if (mask < 2) return;                           // no biology in solid/wall

    const double S = std::max(C[0], 0.0);

    for (std::size_t m = 0; m < B.size(); ++m) {
        const double Bm = std::max(B[m], 0.0);
        if (Bm <= 0.0) continue;

        const double growth = mu_max[m] * S / (Ks[m] + S) * Bm;   // gDW/L/s

        bioR[m]  += growth - kd[m] * Bm;
        subsR[0] -= growth / Y[m];
        subsR[1] += growth * fP[m];
    }
}
\end{lstlisting}

\section{The six rules}

\begin{enumerate}[leftmargin=1.4em]
\item \textbf{Indices follow \file{CompLaB.xml}.} \code{C[0]} is
\tg{substrate0}; \code{B[0]} is \tg{microbe0}. There is no lookup by name, so
reordering substrates in the XML silently changes what your rate law means.
Keep a comment naming each index at the top of the file.

\item \textbf{Everything is per second.} A rate constant published per day must
be divided by 86400.

\begin{cbwarn}
\textbf{\color{cbRed}This causes more dead runs than anything else in the
code.}\quad A rate 86400 times too large blows up on the first step; a rate
86400 times too small does nothing at all and looks like a configuration
problem. Write the unit in a comment next to every constant.
\end{cbwarn}

\item \textbf{Zero the output arrays first.} They are reused between voxels. A
rate law that returns early without zeroing leaks the previous voxel's answer
into this one.

\item \textbf{Guard on size, and match your chemistry.} The
\code{if (C.size() < 2)} line is not decoration. The shipped uranium network
guards with \code{C.size() < 14}; drop that file into a two-substrate case and
it reads past the end of the array. \file{validateExamples.py} checks this
guard against the substrate count for exactly that reason.

\item \textbf{Floor negatives.} \code{std::max(x, 0.0)} on every concentration
you read. The lattice-Boltzmann solver can produce small negative values near
steep fronts, and a rate law that squares one amplifies noise into an
instability.

\item \textbf{Never call \code{exit()}.} Under MPI it leaves the other ranks
hanging and destroys the run. Print a warning and return a safe value; a voxel
that grows at zero rate is visible in the output and recoverable.
\end{enumerate}

\section{Checking a rate law before you trust it}

Mass balance is the cheapest real test, and you can do it without Palabos.
Compile the header on its own against a stub, feed it random states, and check
the conservation you know must hold:

\begin{lstlisting}[style=cpp]
// mini driver: A + B -> C must satisfy dA + dC = 0 and dB + dC = 0
for (int trial = 0; trial < 20000; ++trial) {
    std::vector<double> C = {rnd(), rnd(), rnd()}, R(3);
    defineAbioticRxnKinetics(C, R, 2);
    assert(std::fabs(R[0] + R[2]) < 1e-12);
    assert(std::fabs(R[1] + R[2]) < 1e-12);
}
\end{lstlisting}

\file{validateExamples.py} does the compile half of this for every example; the
balance half is worth writing for your own chemistry, once.

% =====================================================================
\chapter{Training your own surrogate}
\label{ch:trainsurrogate}

The network shipped in \file{surrogateModel.hh} encodes one organism, two
substrates, one range. For anything else you fit your own. Everything for that
is in \file{surrogate\_training/}.

\section{The three steps}

\begin{center}
\begin{tikzpicture}[font=\footnotesize,
  s/.style={draw=cbRule,fill=cbSand,rounded corners=2pt,minimum width=3.5cm,
            minimum height=0.9cm,align=center,line width=0.7pt},
  a/.style={-{Latex[length=2mm]},cbGrey,line width=0.7pt}]
\node[s] (m) at (0,0)    {your model\\ \texttt{.xml / .mat / .json}};
\node[s] (d) at (0,-1.8) {\texttt{training\_data.csv}};
\node[s] (h) at (0,-3.6) {\texttt{surrogate\_weights\_*.hh}};
\node[s] (v) at (0,-5.4) {verified against the fit};
\draw[a] (m) -- node[right,font=\tiny,xshift=2pt]{\texttt{generateTrainingData.py} \ \ offline FBA sweep} (d);
\draw[a] (d) -- node[right,font=\tiny,xshift=2pt]{\texttt{trainSurrogate.py} \ or \ \texttt{trainSurrogate.m}} (h);
\draw[a] (h) -- node[right,font=\tiny,xshift=2pt]{\texttt{verifyExport.py} \ \ compiles it and compares} (v);
\end{tikzpicture}
\end{center}

\section{Step 1: sweep the FBA offline}

\begin{lstlisting}[style=sh]
python3 generateTrainingData.py iAF987.xml \
    --exchange EX_ac_e   --range 1e-3 10   --log \
    --exchange EX_fe3_e  --range 1e-5 0.5  --log \
    --samples 4000 --seed 1 \
    -o geobacter_training.csv
\end{lstlisting}

\subsection{Choosing the range}

\begin{cbnote}
Use the range the \textbf{simulation} will visit, not the range the organism
can tolerate. CompLB3D's uptake bound is
$V = V_{\max} C/(K_c+C)$, so the largest uptake any voxel can ever request is
$V_{\max}$ --- your \tg{fba\_maximum\_uptake\_flux}. That is the natural upper
end of the sweep. The lower end should be small but non-zero.
\end{cbnote}

Sample \textbf{logarithmically} whenever the range spans more than about one
order of magnitude. Pore-network concentrations routinely span four, and a
linear grid puts nearly all its points in the top decade --- so the fit is worst
exactly where the biology is most sensitive.

Keep to \textbf{1 to 3 inputs}. Beyond that the sample count needed for the
same accuracy grows fast and the fit stops being worth the trouble.

\section{Step 2: fit and export}

Two paths, writing byte-identical output. Use whichever you have:

\begin{center}
\begin{tabular}{@{}p{4.8cm}p{10.0cm}@{}}
\toprule
\file{trainSurrogate.m} & MATLAB, Deep Learning Toolbox, \code{trainlm}. Reconstructs how the shipped network was made.\\
\file{trainSurrogate.py} & numpy and scipy, L-BFGS-B on analytic gradients. No licence needed.\\
\bottomrule
\end{tabular}
\end{center}

\begin{lstlisting}[style=sh]
python3 trainSurrogate.py geobacter_training.csv \
    --name geobacter --microbe-id 0 \
    -o surrogate_weights_geobacter.hh
\end{lstlisting}

Both default to the shipped architecture --- 4 hidden layers of 10 tansig
neurons and a linear output, with \code{mapminmax} scaling on both ends --- so a
retrained network drops into an existing build with no other change.

\section{Step 3: verify}

\begin{lstlisting}[style=sh]
python3 verifyExport.py surrogate_weights_geobacter.hh
\end{lstlisting}

\begin{cbnote}
\textbf{\color{cbTeal}Do not skip this.}\quad An exporter that transposes a
matrix or drops a digit produces a header that compiles, runs, and is wrong,
with no independent reference to notice. So the trainer writes 512 reference
predictions next to the header, and \file{verifyExport.py} compiles the header
and checks them. It should report agreement to about $10^{-13}$ relative.
\end{cbnote}

\section{Installing it}

The generated header is namespaced, so several organisms coexist:

\begin{lstlisting}[style=cpp]
#include "surrogate_weights_geobacter.hh"
...
else if (microbeId == 0) {
    bioR[microbeId] = srgnet_geobacter::eval(Fin[microbeId]);
}
\end{lstlisting}

Then \tg{enable\_surrogate} \code{true} and
\tg{reaction\_type} \code{surrogate} in the XML.

Every generated header also carries \code{inTrainingRange()}. Call it if a run
may wander outside the box.

\section{Inspecting a network somebody else fitted}

\begin{lstlisting}[style=sh]
python3 inspectSurrogate.py surrogateModel.hh
\end{lstlisting}

The training range is not stored anywhere as such --- but \code{mapminmax} is
invertible, so it can be recovered exactly:
\[
  \text{gain} = \frac{2}{x_{\max}-x_{\min}},\quad
  \text{offset} = x_{\min}
  \qquad\Longrightarrow\qquad
  x_{\max} = \text{offset} + \frac{2}{\text{gain}}
\]

Run it before trusting any network, including the shipped one.

% =====================================================================
\chapter{Running on a cluster}
\label{ch:cluster}

\section{MPI}

MPI is on by default. Nothing in the input file changes; you launch with
\code{srun} or \code{mpirun} instead of running the binary directly.

\begin{lstlisting}[style=sh]
srun ./complab
\end{lstlisting}

Palabos decomposes the domain into blocks, one per rank, with a halo. That has
a consequence worth internalising:

\begin{cbwarn}
\textbf{\color{cbRed}More ranks is not always faster.}\quad Each rank needs
enough interior voxels to be worth the halo exchange. The shipped examples are
3456 voxels; on 36 ranks that is about 96 voxels each, nearly all of them halo,
and the run is slower than on one core. Ask for ranks in proportion to your
domain, not to what the queue will give you.
\end{cbwarn}

\section{A Slurm script}

\file{comp.sh} is a production template; \file{examples/runExamples.slurm}
submits the example suite. The parts that matter:

\begin{lstlisting}[style=sh]
#SBATCH --account=PROJECT          # required, or the job is rejected
#SBATCH --nodes=1
#SBATCH --ntasks-per-node=36       # match this to your domain size
#SBATCH --time=04:00:00

cd $SLURM_SUBMIT_DIR               # complab reads CompLaB.xml from the cwd
module purge
module load gcc openmpi            # THE SAME modules you built with
srun ./complab
\end{lstlisting}

\note{Load the same modules you built with.}{A binary built against one MPI and
launched under another fails in ways that look like your code is broken. Put
the \code{module load} line in both the build script and the run script, and
keep them identical.}

\section{Optional dependencies on a cluster}

\begin{center}
\begin{tabular}{@{}p{3.4cm}p{11.4cm}@{}}
\toprule
GLPK & \code{module load glpk}, or build it in your home directory and point \file{CMakeLists.txt} at it\\
COBRApy & needs cobrapy importable \textbf{at run time}, not just at build time, and \file{src/complab3d\_cobrapy.py} present at \tg{src\_path}. A virtualenv that is not activated in the job script is the usual failure.\\
\bottomrule
\end{tabular}
\end{center}

\section{Restarting}

Long runs should checkpoint. Set \tg{save\_CHK\_interval} to something
non-zero, and resume with:

\begin{lstlisting}[style=xml]
<read_NS_file>true</read_NS_file>
<read_ADE_file>true</read_ADE_file>
\end{lstlisting}

The flow field and the chemistry restart independently, which is useful: a
converged flow field on a fixed geometry can be reused across many chemistry
runs.

% =====================================================================
\chapter{Output and post-processing}
\label{ch:output}

\section{What gets written}

\begin{center}
\begin{tabular}{@{}p{4.2cm}p{3.0cm}p{7.4cm}@{}}
\toprule
\textbf{File} & \textbf{Format} & \textbf{Contents}\\
\midrule
\code{nsLattice*.vti} & VTK ImageData & velocity and pressure\\
\code{maskLattice*.vti} & VTK ImageData & material numbers, which change if minerals do\\
\code{subsLattice}$n$\code{*.vti} & VTK ImageData & substrate $n$\\
\code{bioLattice}$n$\code{*.vti} & VTK ImageData & microbe $n$\\
\code{*.chk} & Palabos binary & checkpoints, for restart\\
\bottomrule
\end{tabular}
\end{center}

The index and the iteration number are appended to each filename, so
\file{subsLattice0\_0000500.vti} is substrate 0 at step 500.

\section{ParaView}

\begin{enumerate}[leftmargin=1.4em]
\item \textbf{Open} the numbered group; ParaView collapses the time series into
one entry.
\item \textbf{Apply}, then set \textbf{Coloring} to the field.
\item \textbf{Slice} to see inside. The default view of a 3D block shows only
its outer surface, which for a pore network is almost useless.
\item \textbf{Threshold} on the mask field to hide solid voxels.
\end{enumerate}

\note{Unit spacing.}{The \file{.vti} files are written with spacing 1, so
ParaView shows lattice coordinates. Multiply by \tg{dx} for physical distance,
or set the spacing in ParaView's Information panel.}

\section{Getting numbers out}

For anything quantitative --- a breakthrough curve, a mass balance, a converged
biomass ratio --- read the \file{.vti} files with a script rather than reading
values off a picture.

\begin{lstlisting}[style=py]
# pip install vtk numpy
import vtk
from vtk.util.numpy_support import vtk_to_numpy

r = vtk.vtkXMLImageDataReader()
r.SetFileName("output/subsLattice0_0000500.vti")
r.Update()
img  = r.GetOutput()
dims = img.GetDimensions()
arr  = vtk_to_numpy(img.GetPointData().GetArray(0)).reshape(dims[::-1])

print("total:", arr.sum())          # for mass balance
print("mid-plane profile:", arr[dims[2]//2, dims[1]//2, :])
\end{lstlisting}

\section{The checks worth automating}

\begin{center}
\begin{tabular}{@{}p{4.2cm}p{10.6cm}@{}}
\toprule
\textbf{Check} & \textbf{Why}\\
\midrule
totals over time & mass balance. A conservative tracer's total should change only through the boundaries.\\
$z$-independence & every shipped example is quasi-2D; a $z$ dependence is a bug.\\
negative minimum & should never be below zero by more than round-off.\\
porosity & for precipitation and dissolution runs, the number that tells the story.\\
\bottomrule
\end{tabular}
\end{center}
